% 原书第 227 页（PDF 第 240 页）
\BookPage{227}{240}
\BookHeading{盛清扬州的城市生活—以《扬州画舫录》为中心}{盛清扬州的城市生活\\[3pt]—以《扬州画舫录》为中心}{常建华*}
\begin{BookQuote}
「扬一益二」，扬州自古就是富庶之地。隋唐以来扬州又成为东南地区的繁华
\end{BookQuote}

都市，唐代诗人杜牧有句：「谁知竹西路，歌吹是扬州。」至清代十八世纪，扬州盐业、商业发达，发财致富的徽州商人在扬州经商的同时，进行奢侈消费，加之清帝在扬州南巡，地方官与商人奢华接待，都推动了扬州的城市建设与城市生活的繁荣。十八世纪之后，清朝盛极而衰，南巡停止，盐商日渐中落，扬州也失去了往日的辉煌。十八世纪的扬州，是盛清的表征，也在一定程度上代表着中国古代城市的发展水准。

清代仪征人李斗，字艾塘，又字北有，喜好搜集扬州掌故，所着《扬州画舫录》成书于乾隆末期，记录了盛清扬州的繁华。学人对于该书非常重视，评价甚高。专门研究笔记的谢国桢先生认为：该书「概括而有系统地叙述了扬州都市与社会风俗的情况」。来新夏先生指出：「画舫录本为记士大夫消闲活动之作。《扬州画舫录》虽用此书名，但所记内容已不仅限于歌场舞榭，乃记扬州地方风土人情、人文胜迹、社会风貌的地方史志。」²在扬州地方文史学者韦明铧先生看来，《扬州画舫录》「是用普通人的眼睛去看一座城市」，从所记载的日常小事，使我们「深深地感受到一个太平的世道和一切真切的生活」，「开辟了对于繁华社会生活的另类写法」3。可以说，《扬州画舫录》是认识清代城市生活最好的资料之一。

\vfill
\begin{BookNotes}
\BookNote{*}{作者任职于南开大学中国社会史研究中心。}
\BookNote{1.}{谢国桢，《明清笔记谈丛》，上海:上海古籍出版社，1981年版，第144页。}
\BookNote{2.}{来新夏，《清人笔记随录》，北京：中华书局，2005年版，第264页。}
\BookNote{3.}{韦明铧，《扬州掌故》，苏州：苏州大学出版社，2001年版，第92、93、94页。}
\end{BookNotes}

% 原书第 228 页（PDF 第 241 页）
\BookPage{228}{241}
近十几年来，海内外探讨扬州城市社会文化的著述不断出现4，其中也有城市生活的论述。不过以往的研究比较缺乏「日常生活」的角度，所论问题也还有余地。虽然学者研究扬州都利用过《扬州画舫录》，但是由于该书内容「杂七杂八」的丰富性，有关城市生活的诸多细节还可以深入挖掘。因此，笔者不揣谢陋，以《扬州画舫录》为基本资料，权作续貂之作，以为练笔。

\BookSubheading{一、清前期的扬州城市空间}

地处长江与运河交会处的扬州，不仅是重要的交通孔道，也是盐业产地，促进了城市的发展。明清时期扬州有新城、旧城之分，旧城，即宋大城之西南隅，元至正十七年（1357）朱元璋军队攻克扬州，佥院张德林始改筑之，约十里，周围一千七百七十五丈五尺，高倍之。门五，曰海宁，清改称大东，曰通泗，清改称西门，曰安江，清改称南门，曰镇淮，请改称北门，曰小东。南北水关二，引市河水以通于濠。随着明代中期以后盐政制度改革与商品经济的发展，大量山陕、徽州商人来到扬州，居住在旧城东面临近运河一带。

新城，即宋大城之东南隅，明嘉靖三十四年（1555），倭寇侵掠扬州，为保护盐课收入与盐商安全，知府吴桂芳、石茂华先后筑城。自旧城东南角起，折而南，循运河而东，折而北，复折而西，极于旧城之东北角止。东与南北三面，约八里有奇，计一千五百四十二丈。门七，曰挹江，清改称钞关；曰便门，清改称徐宁；曰拱宸，清改称天宁；曰广储；曰便门，清改称便益；曰通济，清改称缺口；曰利津，清改称东关。沿旧城城濠南北水关二，东与南二面，即以运河为城濠，北面作濠，与旧城连，注于运河。（卷9，第178页）

增筑新城，城市空间大为扩展。新城又称东城、石城，发展迅速，至万历年

\vfill
\begin{BookNotes}
\BookNote{4.}{国内的有关论述如冯尔康等，《扬州研究：江都陈轶群百龄冥诞纪念论文集》，台北：联经出版公司，1996年；王振忠，《明清徽商与淮扬社会变迁》，北京：三联书店，1996年；刘凤云，《明清城市空间的文化探析》，北京：中央民族大学出版社，2001年。国外的研究如美国学者梅尔清《清初扬州文化》，上海：复旦大学出版社，2004年；澳大利亚学者安东篱，《说扬州：1550-1850年的一座中国城市》，北京：中华书局，2007年。}
\BookNote{5.}{本文使用的《扬州画舫录》是江苏广陵古籍刻印社1984年版，由周光培点校。为从简，本文引自该书采取夹注，注明卷数与页数。}
\end{BookNotes}

% 原书第 229 页（PDF 第 242 页）
\BookPage{229}{242}
\noindent 间，形成处新城者尽富商大贾，居旧城多缙绅之家的社会空间格局。

清初扬州遭受清军毁坏，居民死伤惨重。随着康熙中后期社会安定与经济的恢复，扬州盐业得到发展。扬州城西北的平山堂因是宋代欧阳修任扬州知府时流连往返的名胜之地，成为观光地。下为蜀冈，保障河流经其上。雍正十年（1723)知府尹会一募捐重浚保障河，此后乾隆十五年、二十年、二十六年巡盐御史吉庆、普福、高恒「叠经疏浚，加深广曲折，点缀园亭，栽植桃柳」。需要指出，乾隆十六年、二十年、二十七年是皇帝前三次南巡的年份，显然乾隆时期疏浚保障河工程是为了皇帝巡视扬州。而为了疏浚工程的「募捐」则来自于盐商。乾隆二十三年(1758）又凿莲花埂，河通平山堂。此后，保障河畔成为著名的冶游之地，人称「瘦西湖」。

旧城市河连接城内外，也得到了疏浚，沟通新城，以供新城盐商富户冶游。从新城市河到西北保障河湖，游船往来，成为专为旅游设置的城内市河，有「小秦淮」之称。

为了乾隆帝南巡驻跸扬州，从城外东北的香阜寺抵达新城天宁寺行宫，两淮商人开设「新河」，便于乾隆帝巡游。由天宁门、天宁街、彩衣街、司前三铺、教场、辕门桥、多子街、埂子上，出钞关、花觉行，至钞关马头为新城御道。道旁或搭彩棚，或陈水嬉。

新城商业区迅速拓展。乾隆三十二年（1767）总督高晋将原新城卫教场改移西门外蜀冈之下，原有空地除供营兵操练外，余者由盐商召民领建认租，作为盐务公地。由天宁门城内东入彩衣街，左折运司街、教场、辕门桥、多子街、埂子上，出钞关门，右折花觉行，入九峰园，此小东门外新城御道。「南巡时，墁石清道，如铺沙藉路之例。」（卷9，第185页）御道也成为扬州商业最为繁华的地段7。盐商多居新城内南、北河下，相近的丁家湾，凡替商家经手者，俱集于此，有些人租几

\vfill
\begin{BookNotes}
\BookNote{6.}{（清）赵之璧着，《平山堂图志》卷1〈名胜上·炮山河〉，扬州：广陵书社，2004年版，第11页。}
\BookNote{7.}{有关扬州城区不同功能区的划分，参考了王振忠《明清徽商与淮扬社会变迁》第74-87页的有关论述。}
\end{BookNotes}

% 原书第 230 页（PDF 第 243 页）
\BookPage{230}{243}
\noindent 间房子，名曰「公店」，「任买卖人往来交易」8。

城外扬州名胜以园林著称。「扬州园林之胜，甲于天下。由于乾隆朝六次南巡，各盐商穷极物力以供宸赏，计自北门直抵平山，两岸数十里楼台相接，无一处重复。其尤妙者在虹桥迤西一转，小金山矗其南，五顶桥锁其中，而白塔一区雄伟古朴，往往夕阳返照，箫鼓灯船，如入汉宫图画。」盐商以重资广延名士为之创稿，一一布置，规划了城外风景区。扬州名园多为盐商所建，城外分布于保障河两岸，城内园林区主要分布新城，以南河下盐商居住区为多。「城内之园数十，最旷逸者，断推康山草堂。而尉氏之园，湖石亦最胜，闻移植时费二十余万金。其华丽缜密者，为张氏观察所居，俗所谓张大麻子是也。」10张氏之园广数十亩，中有三层楼，可瞰大江，凡赏梅、赏荷、赏桂、赏菊，皆各有专地。演剧宴客，上下数级如大内式。另有套房三十余间，回环曲折，迷不知所向。金玉锦绣，四壁皆满，禽鱼尤多。

清代扬州画舫皆在城外，惟大东、小东二门码头在城中。

江都、甘泉二县同附郭，旧城西半壁，新城南半壁，为江都治；旧城东半壁，新城北半壁，为甘泉治。

关于扬州的街道。旧城，南门至北大街，三里半，近南门者谓之南门大街，近北门者谓之北门大街，中谓之院大街。大东门至西门，一里半，近大东门者谓之大东门大街，近西门者谓之西门大街。南水关抵北水关市河，东岸直街为南门、盐院、北门三大街，西岸直街为南、中、北三小街。小东门城脚至大东门城脚有九条巷，其上有两层街。

新城，东关至大东门大街，三里，近东关者谓之东关大街，近大东门者谓之彩衣街。阙口门至小东门大街，三里，近阙口门者谓之阙口大街，上为左卫街、多

\vfill
\begin{BookNotes}
\BookNote{8.}{（清）林苏门撰，《邗江三百吟》卷1〈丁家湾公店〉引言，该书嘉庆十三年成书，广陵书社，2005年版，第15页。}
\BookNote{9.}{（清）欧阳兆熊、金安清，《水窗春》卷下〈维扬胜地〉，北京：中华书局，1984年版，第72页。}
\BookNote{10.}{（清）欧阳兆熊、金安清，《水窗春艺》卷下〈维扬胜地〉，中华书局，1984年版，第72页。}
\end{BookNotes}

% 原书第 231 页（PDF 第 244 页）
\BookPage{231}{244}
\noindent 子街，抵小东门街。钞关至天宁门大街，三里半，近钞关者谓之埂子上，上为南柳巷、北柳巷，至天宁门，谓之天宁门大街。钞关东沿内城脚至东关，为河下街。自钞关至徐宁门，为南河下，徐宁门至阙口门，为中河下；阙口门至东关，为北河下，计四里。徐宁门至罗湾止，计二里，由徐宁门至蒋家桥，为徐宁门大街，由蒋家桥至罗湾，为皮市街。湾子上为城中斜街，自罗湾起至打铜巷止。

城河即市河，南出龙头关，有坝蓄水，与官河隔，谓之针桥；北出大东门水关至高桥，亦有坝蓄水，与官河隔，谓之黄金坝，此古市河。清龙头关淤垫，乃于小东门钓桥下筑坝，令河北徙，出大东门水关，汇镇淮门市河，入保障湖，以利东城画舫，凡小东门外城脚头敌台、二敌台、头巷、二巷，皆画舫码头。

扬州商业性街市发展明显，成为城市繁荣的重要空间。如天宁门至北门，沿河北岸建河房，仿京师长连短连、廊下房及前门荷包棚、帽子棚做法，谓之「买卖街」。「令各方商辇运买珍异，随营为市。」（卷4，第100页）

\BookSubheading{二、园林与城市生活}

遍游江浙的士大夫刘大观说：「杭州以湖山胜，苏州以市肆胜，扬州以园亭胜，三者鼎峙，不可轩轻。」（卷6，第144页）可见扬州的园林成为城市标志性景观。金安清将扬州比较南京、苏州、杭州后说：「扬州则全以园林亭榭擅场，虽皆由人工，而匠心灵构，城北七八里夹岸楼舫无一同者，非乾隆六十年物力人才所萃，未易办也。」11江南人钱泳也说到：「扬州之平山堂，余于乾隆五十二年秋始到，其时九峰园、倚虹园、筱园、西园曲水、小金山、尺五楼诸处，自天宁门外起直到淮南第一观，楼台掩映，朱碧鲜新，宛入赵千里仙山楼阁中。」12当时园林之盛，跃然纸上。

扬州的园林，往往也称为园亭。由香阜寺一站至塔湾，其蜀冈三峰及黄、江、程、洪、张、汪、周、王、闵、吴、徐、鲍、田、郑、巴、余、罗、尉诸园亭。李斗记载，扬州:「郡城以园胜。康熙间有八家花园，王洗马园即今舍利庵，卞园、

\vfill
\begin{BookNotes}
\BookNote{11.}{（清）欧阳兆熊金安清，《水窗春》卷下〈广陵名胜〉，中华书局，1984年版，第46页。}
\BookNote{12.}{（清）钱泳，《履园丛话》卷20〈园林·平山堂〉，中华书局，1979年版，下册，第533页。}
\end{BookNotes}

% 原书第 232 页（PDF 第 245 页）
\BookPage{232}{245}
\noindent 员园在今小金山后方家园田内，贺园即今莲性寺东园，冶春园即今冶春诗社，南园即今九峰园，郑御史园即今影园，条园即今三贤祠。」（卷1，第22页）

扬州园林，种植花卉，盆景享有盛名。李斗有详细的记载，湖上园亭，皆有花园，为莳花之地。桃花庵花园在大门大殿阶下。养花人谓之花匠，莳养盆景，蓄短松、矮杨、杉、柏、梅、柳之属。海桐、黄杨、虎刺以小为最，花则月季、丛菊为最，冬于暖室烘出芍药、牡丹，以备正月园亭之用。盆以景德窑、宜兴土、高资石为上等。种树多寄生，剪丫除肄，根枝盘曲而有环抱之势。其下养苔如针，点以小石，谓之「花树点景」。「又江南石工以高资盆增土叠小山数寸，多黄石、宣石、太湖、灵璧之属，有孔有岫，有罅有杠，蓄水作小瀑布倾泻危溜。其下空处有沼，畜小鱼游泳嚅，谓之『山水点景』。」（卷2，第34页）

园林也是聚会之所，士大夫通过诗会雅集，饮酒赏曲，丰富文化生活。李斗说，扬州诗文之会，以马氏小玲珑山馆、程氏筱园及郑氏休园为最盛。至会期，于园中各设一案上置笔二、墨一、端砚一、水注一、笺纸四、诗韵一、茶壶广、碗一、果盒茶食盒各一，诗成即发刻，三日内尚可改易重刻，出日遍送城中矣。每会酒肴俱极珍美，一日共诗成矣。请听曲，邀至一厅甚旧，有绿琉璃四，又选老乐工四人至，均没齿秃发，约八九十岁矣，各奏一曲而退。倏忽间命启屏门，门启则后二进皆楼，红灯千盏，男女乐各一部，俱十五六岁妙年也。李斗还听说，诗牌以象牙为之，方半寸，每人分得数十字或百余字，凑集成诗，最难工妙。休园、筱园最盛。（卷8，第172页）关于马氏小玲珑山馆，有看山楼、红药阶、七峰草堂、清响阁、藤花书屋、丛书楼、觅句廊、浇药井、梅寮诸胜，马氏日管、日璐兄弟俱能诗，好客。

扬州的园林很多。根据钱泳《履园丛话》卷二十〈园林〉的记载，可以举出如江园，扬州江畹香侍郎家所有，在阙口门大街，回廊曲榭，花柳池台，直可与康山争胜。静修俭养之轩，在天宁门内，鲍肯园赠公所筑。四围楼阁，通以廊庑，阶前湖石数峰，尽栽丛桂、绣球、丁香、白皮松之属。九峰园，奇石玲珑，其最高者有九，故以名园。

园林的开放性与旅游。湖边的园林可以供人参观。虹桥附近有盐商洪氏的园


% 原书第 233 页（PDF 第 246 页）
\BookPage{233}{246}
\noindent 林，其中的妙远堂，就是园中待游客地。李斗说:「湖上每一园必作深堂，饬庖寝以供岁时宴游，如是堂之类。」（卷10，第210页）他还介绍，致佳楼楼后皆新辟荒地，并转角桥西口之冶春茶社围入园中。园三面临水，中筑桂花书屋，逶迤连络小室数十间，「令游者惝恍弗知所之」。（卷10，第210页）

园林是文人的理想栖息地。江宁人熊之勋，工诗善书。其家有「小西湖」之胜，然而常居扬州江春的康山草堂。江春表甥林道源工诗，不存稿，居是园十年。

三·食衣住行

扬州人的生活消费依赖市场，颇为时尚。

1.茶酒食肆遍布，讲究饮食

\BookSubheading{(1) 茶肆}

扬州饮茶之所名为茶肆，著称于世。李斗说:「吾乡茶肆，甲于天下，多有以此为业者。出金建造花园，或鬻故家大宅废园为之。楼台亭舍，花木竹石，杯盘匙箸，无不精美。」（卷1，第25页）茶肆分为荤、素两种类型，城中荤茶肆之最盛者，辕门桥有二梅轩、蕙芳轩、集芳轩，教场有腕腋生香、文兰天香，埂子上有丰乐园，小东门有品陆轩，广储门有雨莲，琼花观巷有文杏园，万家园有四宜轩，花园巷有小方壶；素茶肆之最盛者，则有天宁门之天福居，西门之绿天居。茶肆中，城外占湖山之胜，北门桥的双虹楼为最。楼五楹，东壁开牖临河，可以眺远。

茶肆的点心各具特色。双虹楼烧饼，开风气之先，有糖馅、肉馅、干菜馅、苋莱馅之分；宜兴丁四官开蕙芳、集芳，以糟窖馒头得名，二梅轩以灌汤包子得名，雨莲以春饼得名，文杏园以稍麦得名，谓之鬼蓬头，品陆轩以淮饺得名，小方壶以菜饺得名，各极其盛。「而城内外小茶肆或为油镞饼，或为甑儿糕，或为松毛包子，茆檐华门，每旦络绎不绝。」（卷1，第26页)

茶肆颇密集。史阁部墓侧的柳林，有屋十数间为茶肆，题其帘曰「柳林茶社」。（卷3，第76页）


% 原书第 234 页（PDF 第 247 页）
\BookPage{234}{247}
岁时节日，光顾茶肆者众多，生意最好。城北的小洪园后门为旧时「且停车茶肆」，其旁为「七贤居」，也是茶肆。二肆最盛于清明节放纸鸢、端午龙船市、九月重阳九皇会，斗蟋蟀，看菊花，是岁时记中胜地」（卷6，第140页）。

明月楼茶肆在二钓桥南，南岸外为二道沟，中皆淮水，逢潮汐则江水间之。肆中茶取于是，饮者往来不绝，人声喧阗，杂以笼养鸟声，隔席相语，恒以眼为耳。（卷7，第158页）

茶肆颇为热闹。城南的明月楼茶肆在二钓桥南，南岸外为二道沟，中皆淮水，逢潮汐则江水间之。「肆中茶取于是，饮者往来不绝，人声喧阗，杂以笼养鸟声，隔席相语，恒以眼为耳。」（卷7，第158页）到了「以眼为耳」的地步，可见人声鼎沸。

青莲斋在新城天宁街街西，为六安山僧茶叶馆。僧有茶田，春夏入山，秋冬居肆，东城游人，皆于此买茶供一日之用。

\BookSubheading{(2) 酒肆}

酒肆也很盛。北郊酒肆，自醉白园始，康熙间如野园、冶春社、七贤居、且停车之类，皆在虹桥。壶觞有限，不过游人小酌而已。后里人韩醉白于莲花埂构小山亭，游人多于其家聚饮，因呼之曰韩园。迨醉白死，北门街构食肆慕其名而书之，谓之「醉白园」。「园之后门，居小迎恩河西岸，画舫多因之饮食焉。」

新城天宁街的街西，有扑缸春酒肆。李斗说：「游屐入城，山色湖光，带于眉宇，烹鱼煮笋，尽饮纵谈，率在于是。」（卷4，第78页）

城南砚池北有秀野园酒肆，周围景色宜人，「春暖莺飞，禽声杂出，湖外黄花烂缦，千顷一色」。（卷7，第159页）

大仪乡蜀冈有观音街，中建过街亭，凡上山香客于此憩止。「亭外多酒肆，以郭章为最，名得胜园。」（卷16，第347页）

扬州买酒的商店，各有经营特色。李斗说:「扬州市酒以戴氏为最，谓之戴


% 原书第 235 页（PDF 第 248 页）
\BookPage{235}{248}
\noindent 蛮，次则周氏，谓之周六槽坊，皆鬻木瓜酒，若镇江府百花酒，扬州盛行之，则有郭咸泰。」（卷4，第177页）

小东门西城下间有星货铺，即散酒店、庵酒店之类。（卷9，第188页）

虹桥西岸有跨虹阁，为酒铺。阁外日揭帘，夜悬灯，帘以青白布数幅为之，下端裁为燕尾，上端夹板灯，上贴一「酒」字。土酒如通州雪酒、泰州枯、陈老枯、高邮木瓜、五加皮、宝应乔家白，皆为名品，而游人则以木瓜为重。李斗说：「近年好饮绍兴，间用百花；今则大概饮高粱烧，较本地所酿为俗矣。」（卷13，第279页）

李斗介绍了造酒过程与种类:造酒家以六月三伏时造曲，曲有米麦二种，受之以范，其方若砖。立冬后煮瓜米和曲，谓之起酵，酒成谓之「醅酒」。瓜米者，糯稻碾五次之称。碾九次为茶米，用以作糕粽；五六次者为瓜米，用以作酒，亦称酒米。醅酒即木瓜酒，以此米可造木瓜酒，故曰瓜米。酒用米曲则甘美，用麦曲则苦烈。烧酒以米为之，曰米烧；以麦为之，曰麦烧。又有自酵酒糟中蒸出，谓之糟烧。其高梁、荞麦、绿豆均可蒸，亦各以其谷名为名。

城外村庄中人善造酒，城内之烧酒，大抵俱来自城外，驴驼车载，络绎不绝。

新酒与旧藏烧酒按照比例相兑，或者烧酒与木瓜酒相兑，产生新的酒。即秋成新刍，为「时酒」，又曰「红梅酒」，一曰「生酒」。时酒一斤，合烧酒之半，曰「火对」。合烧酒十分之二，曰「筛儿」。合挠酒、醅酒各均，为「木三对」。烧酒还有些次生种类。如烧酒未蒸者，为「酒娘儿」，饮之鲜美。以泉水烧酒和之，则成烧蜜酒。

饮酒习俗丰富多彩。其与时令有关，八月「红梅」新熟，各肆择日贴帖，曰开生」，人争买之，曰「尝生」。至二月惊蛰后止，谓之「剪生」。称酒有特色，酒以斤重，斤以笳计，笳以大竹为之，自一两至一斤，笳准此为大小。卖酒有业俗，铺中敛钱者为掌柜，烫酒者为酒把持。凡有沽者斤数，掌柜唱之，把持应之，遥遥赠答，自成作家，殆非局外人所能猝辨。酒铺例为人烫蒲包豆腐干，谓之「旱团鱼」


% 原书第 236 页（PDF 第 249 页）
\BookPage{236}{249}
\BookSubheading{(3)食肆}

扬州的饮食著名。人称「大抵郡人以盐务为生，率习于浮华而精肴馔也」13

城内食肆多附于面馆，面有大连、中碗、重二之分。冬用满汤，谓之「大连」；夏用半汤，谓之「过桥」。面有浇头，以长鱼、鸡、猪为三鲜。大东门有如意馆、席珍，小东门有玉麟、桥园，西门有方鲜、林店，缺口门有杏春楼，三祝庵有黄毛，教场有常楼。李斗又介绍，乾隆初年，徽人于河下街卖松毛包子，名「徽包店」，因仿岩镇街没骨鱼面，以鲭鱼为面，名其店为「合鲭」。仿之者有槐叶楼火腿面。合鲭复改为「玉坡」，遂以鱼面胜。李斗还说：徐宁门问鹤楼以螃蟹胜。而接踵而至者，不惜千金买仕商大宅为之。如涌翠、碧芗泉、槐月楼、双松圃、胜春楼诸肆，楼台亭榭，水石花树，争新斗丽，实他地之所无。「其最甚者，鳇鱼、蟑螯、班鱼、羊肉诸大连，一碗费中人一日之用焉。」（卷11，第255页）

城南头桥上有知己食，食肆中宰夫杨氏，工宰肉，得炙肉之法，谓之「熏烧」，肆额「丝竹何如」。人皆不得其解，或以虽无丝竹管弦之盛语解之，谓其意在觞咏。或以丝不如竹，竹不如肉语解之，谓其意在于肉。李斗认为:「市井屠沽，每藉联扁新异，足以致金，是皆可以不解解之也。」（卷7，第158页）可见食肆为了招徕客人，在匾额上下功夫，附庸风雅。

旧城东有小东门。「小东门街多食肆，有熟羊肉店，前屋临桥，后为河房，其下为小东门码头。就食者鸡鸣而起，茸裘氊帽，耸肩扑鼻，雪往霜来，窥食蹊，探庋阁，以金啖庖丁，迟之又久，先以羊杂碎饲客，谓之小吃；然后进羊肉羹饭，人一碗，食余重汇，谓之走锅；漉去浮油，谓之剪尾。狃以成习，亦觉此嚼不恶，惟不能与贪眠者会食。一失其时，残杯冷炙，绝无风味。」（卷9，第188页）小东门西外城脚无市铺，卯饮申饭，半取资于小东门街食肆，多糊炒田鸡、酒醋蹄、红白油鸡鸭、炸虾、板鸭、五香野鸭、鸡鸭杂、火腿片之属，骨董汤更一时称便。

大东门附近在钓桥外，有申申如食肆，是个素食肆。旁有羊肉店，名曰「回回馆」。（卷9，第198页）可见这里的商家，服务物件主要是回民与素食者。

\vfill
\begin{BookNotes}
\BookNote{13.}{（清）林溥撰，《扬州西山小志》之《市肆十一首》，广陵书社，2005年版，第42页。}
\end{BookNotes}

% 原书第 237 页（PDF 第 250 页）
\BookPage{237}{250}
扬州人饮食喜欢海鲜，水产品丰富，人们购置新鲜鱼蟹消费。府城西北的黄金坝，皆鱼市。扬州城位于江、淮之间，南则三江营，出鳄鱼，瓜洲深港出刀鱼，北则艾陵、甓社、邵伯诸湖，产鱼尤多。由官河乘风而下，城肆贩户于此交易。肆中一日三市，早挑、中挑、晚挑，皆沿湖诸村镇中人为之。李斗说：「村镇设行，渔户取鱼自行交易，挑者输于城中，其行若飞。或三四十里，多至六七十里，不多时即至，以行之迟速分优劣。」（卷1，第16页）蝙鱼、白鱼、鲫鱼为上，鲤鱼、季花鱼、青鱼、黑鱼次之，蝙鱼、罗汉鱼为下，其苍蝙，勒鱼、红蓼鱼、鞋底鱼，则为海鱼。蟹自湖至者为湖蟹，自淮至者为淮蟹，淮蟹大而味淡，湖蟹小而味厚，品蟹者以湖蟹为胜。坝上设八鲜行，八鲜，即菱、藕、芋、柿、虾、蟹、蟑螯、萝卜。此外，城内另有鱼行，鱼市亦谓之鱼摊。

黄金坝鱼行有咸货，有腌切。扬州地居海滨，盐为特产，以盐渍鱼，糗干载入郡城，谓之腌腊。鱼行货物种类繁多，委积尘封，黄养如宁波，海鲤如武昌。最大的鱼是鲨，皮有珠文，肥甘可食。小的鱼以竹贯穿，晾干，最小的为银鱼。沿海渔民拾蛏，鲜者鲍之，不能鲍者干之，其肥在鼻。还有鱼翅，割海鱼鳍制成。海蜇的吃法，已经分成蛰头与蛰皮。石首鱼春产于江，秋产于海，狼山以下人家，八月顿顿食黄鱼。加工石首鱼地方法，「风干其鲷曰膘，取𫗦曰赎，以盐冰之曰腌鱼子，凡此皆行货也」。（卷1，第16页）行货大半属于南货，经营南货的多镇江人，京师称其为「南酒」，所贩皆大江以南之产，又署其店肆曰「海味」。

士大夫讲究饮食。天宁寺西园下院为枝上村，有弹指阁，僧人文思居之。文思工诗，一时乡贤寓公皆与之交朋友。文思善为豆腐羹、甜浆粥，有效法者，谓之「文思豆腐」。（卷4，第83页）

面馆、茶肆也是人际交往的场所。扬州风俗：「交情须会，有不必宴饮者。天明而起，亲至其家，邀往好面馆、好茶肆一餐，名曰『过蚤』。约人议事亦然。」14另外，寻常待客是「三碗六盘」，讲究口味出奇；宴客的丰盛之席为「七簋两点。15

\vfill
\begin{BookNotes}
\BookNote{14.}{（清）林苏门撰，《邗江三百吟》卷9《过蚤》引言，广陵书社，2005年版，第119页。}
\BookNote{15.}{（清）林苏门撰，《邗江三百吟》卷9〈三碗六盘〉、〈七簋两点〉引言，广陵书社，2005年版，第119、120页。}
\end{BookNotes}

% 原书第 238 页（PDF 第 251 页）
\BookPage{238}{251}
2.服饰华丽，引领时尚

迎恩河西，有「邗上农桑」、「杏花村舍」二景。仿康熙帝《耕织图》做法，封隈为岸，建仓房、食盍饷桥、报丰祠。祠前击鼓吹豳台，左有袭房，右有浴蚕房、箔房、绿叶亭。亭外桑阴郁郁，时闻斧声。树间建大起楼，楼下长廊至染色房、练丝房。房外为练池，池外有春及堂。堂右有嫘祖祠、经丝房、听机楼。楼后有东织房、纺丝房。房外板桥二三折，至西织房、成衣房，接献功楼。这一景观浓缩了农业社会的耕织生活。

扬州的服装业发达。有染布业仅次于苏州，李斗说：「江南染房，盛于苏州。扬州染色，以小东门街戴家为最。」（卷1，第28页）

嫘祖祠，祀马头娘。祠右沼堤种竹，竹后长廊数丈，廊竟，横置小舍三间，为经丝房，经机所持丝也。屋右接听机楼。楼台疏处为纺丝房，与经丝房对，居其右。织房十余间，以东西分。「成衣房十余间，纺砖刀尺，声声相闻。」（卷1，第30页)

扬州的服装店繁多，穿衣引领时尚。位于新城御道上的多子街即缎子街，两畔皆缎铺。李斗说：「扬郡着衣，尚为新样。十数年前，缎用八团，后变为大洋莲、拱璧兰；颜色在前尚三蓝、朱、墨、库灰、泥金黄，近用膏粱红、樱桃红，谓之福色。」（卷9，第185页）据说因大将军福康安征讨台湾时过扬州着此色，民间仿效，成为时尚。每当货至，先归绸庄缎行，然后发铺，谓之「抄号」。每年以四月二十日为例，谓之「镇江会」。相当于今天我们说的商品交易会。

扬州服饰的时尚除了颜色之外，样式别出心裁，引领潮流。翠花街，皆珠翠首饰铺。扬州松勒，异于他地，有蝴蝶、望月、花蓝、折项、罗汉松、懒梳头、双飞燕、到枕松、八面观音各种假发1，还有貂覆额、渔婆勒子诸式。穿鞋颇为讲究，女鞋以香樟木为高底，在外为外高底，有杏叶、莲子、荷花诸式；在里者为里高底，谓之道士冠，平底谓之底儿香。无论里高、还是外高，都是通过增加鞋底提高女性的身高。女衫以二尺八寸为长，袖广尺二，外护袖以锦绣镶之；冬则用貂狐

\vfill
\begin{BookNotes}
\BookNote{16.}{（清）林苏门撰，《邗江三百吟》卷6〈松架〉引言记载：「妇女以假发为架，上以真发挽之成松，十年前未曾见此。」可以参考，广陵书社，2005年版，第86页。}
\end{BookNotes}

% 原书第 239 页（PDF 第 252 页）
\BookPage{239}{252}
\noindent 之类。裙式以缎裁剪作条，每条绣花两畔，镶以金线，碎逗成裙，谓之凤尾。这种凤尾穿上，想必金光闪闪，有富贵气。李斗说:「近则以整缎折以细系道，谓之百折；其二十四折者为玉裙，恒服也。」（卷9，第186页）这种折裙17，开始流行。这条街上，还有硝皮袄的「毛毛匠」，加工、制作皮衣，看来皮衣也是时尚之服。

腰间佩带三针表是扬州的时尚。这一习俗受到西洋文化的影响，「洋人一针表，售出价不过数十金；近日面上三针，校定更准，其价更昂。扬城趋时人借定时名目，亦多于腰间佩带，以为饰观」18。

3.居室陈设讲究，具有文化气息

扬州建造房屋非常讲究。钱泳说：「造屋之工，当以扬州为第一，如作文之有变换，无雷同，虽数间小筑，必使门窗轩豁，曲折得宜，此苏、杭工匠断断不能也。盖厅堂要整齐如台阁气象，书房密室要参错如园亭布置，兼而有之，方称妙手。今苏、杭庸工皆不知此义，惟将砖瓦木料搭成空架子，千篇一律，既不明相题立局，亦不知随方逐圆，但以涂汰作生涯，雕花为能事，虽经主人指示，日日叫呼，而工匠自有一种老笔主意，总不能得心应手者也。」19

李斗说，民间厅事，置长几，上列二物，如铜磁器及玻璃镜、大理石插牌，两旁亦多置长几，谓之靠山摆。今各园长几，多置三物，如京式。屏间悬古人画，小室中用天香小几，画案书架。小几有方、圆、三角、六角、八角、曲尺、如意、海棠花诸式。画案长者不过三尺。书架下椟上空，多置隔间。几上多古砚、玉尺、玉如意、古人字画、卷子、聚头扇、古骨朵、剔红蔗葭、蒸饼；河西三撞两撞漆合、磁水盂，极尽窑色，体质丰厚；灵壁、太湖诸砚山、珊瑚笔格；宋蜡笺，书籍皆宋元精椠本、旧抄秘种及毛钞、钱钞。隔间多杂以铜、磁、汉玉古器。他如雉尾

\vfill
\begin{BookNotes}
\BookNote{17.}{（清）林苏门撰，《邗江三百吟》卷6〈百褶裙〉引言记载：「裙用好绸，绉捏成百褶，颜色不拘，利其软而下垂，此新式也。」可以佐证，广陵书社，2005年版，第82页。}
\BookNote{18.}{（清）林苏门撰，《邗江三百吟》卷7〈带三针表〉引言，广陵书社，2005年版，第95页。}
\BookNote{19.}{（清）钱泳《履园丛话》卷12〈艺能·营造〉，中华书局，1979年版，上册，第326-327页。}
\end{BookNotes}

% 原书第 240 页（PDF 第 253 页）
\BookPage{240}{253}
\noindent 扇、自鸣钟2°、螺甸器、银累条、铜龟鹤、日圭、嘉量、屏风合匝、天然木几座、大小方圆古镜、异石奇峰、湖湘文竹、天然木拄杖、宣铜炉。大者为宫奁，皆炭色红、胡桃纹、鹧鸪色，光彩陆离；上品香顶撞、玉如意，凡此皆陈设也。（卷17，第401-402页)

庭院的陈设也很讲究，有鱼缸养观赏鱼以及栽种鲜花。李斗说，城中富家以花事为陈设，更替以时。柳下置砂缸蓄鱼，有文鱼、蛋鱼、睡鱼、蝴蝶鱼、水晶鱼诸类。

4.出行

轿子的使用较为普遍。李斗说，山轿二乘，即竹兜子，闲时贮一粟庵，遇官舟抵岸则出。至四城堂客，上山多步行，富贵家则自备女舆，行走若飞，谓之「飞轿」；步碎而软，谓之「溜步」。轿夫谓之「楼儿」，随轿侍儿谓之「跑楼儿」。(卷16，第350页)

扬州人生活方式中引人瞩目的还有泡澡习尚。李斗说浴池之风，开始于城北邵伯镇的郭堂，后城内徐宁门外之张堂仿效，城内张氏复于兴教寺效其制以相竞尚，于是四城内外皆然。著名的城内有开明桥之小蓬莱，太平桥之白玉池，缺口门之螺丝结顶，徐宁门之陶堂，广储门之白沙泉，埂子上之小山园，北河下之清缨泉，东关之广陵涛。而城外则坛巷之顾堂，北门街之新丰泉最为著名。浴池十分讲究，以白石为池，方丈余，间为大小数格，其大者近镬水热，为大池，次者为中池，小而水不甚热者为娃娃池。贮衣之柜，环列于厅的为座箱，在两旁的为站箱。内通小室，谓之暖房。还有按摩，「茶香酒碧之余，侍者折枝按摩，备极豪侈」。（卷1，第25页），另据记载，扬州澡堂又称「混堂」，「城内外数以百计。凡堂外有立箱、有坐箱、有凉池、有暖房、有茶汤处、有剃头修脚处。堂内之池取乎洁，用白矾石界为三四池，水之温凉，各池不同」21。当地习俗，男子亲迎前一夕入浴，动费数十金。除夕浴谓之「洗邋遢」，端午谓之「百草水」。

\vfill
\begin{BookNotes}
\BookNote{20.}{（清）林苏门撰，《邗江三百吟》卷7〈厅上闹钟〉引言记载：「自鸣钟以定时刻，扬城趋时之家，间亦用之。近日有用闹钟设于厅上者…厅上乃宾客往来之地，借以骇人见闻。」说明扬州开始流行厅中摆设自鸣钟。广陵书社，2005年版，第94页。}
\BookNote{21.}{（清）林苏门撰，《邗江三百吟》卷3〈混堂〉引言，广陵书社，2005年版，第42页。}
\end{BookNotes}

% 原书第 241 页（PDF 第 254 页）
\BookPage{241}{254}
扬州人有早上喝茶、晚上洗澡之俗，人称早上「皮包水」，晚上「水包皮」，反映了闲适、惬意的生活节奏。「这种生活只有都市人能有，农村人自然做不到。」22

\BookSubheading{四、画舫与游乐}

画舫者，即游船也。李斗以画舫为书名，记录扬州，强调的是以游为特色的城市风貌。而游船所经在于河与湖，游湖最能体现扬州的城市生活。李斗指出了画舫的由来：

\begin{BookQuote}
扬州画舫，始于鼓棚。鼓棚本泰州驳盐船，至朽腐不能装载，辄牵入内河，架以枋楣橡柱。大者可置三席，谓之「大三张」，小者谓之「小三张」。驳盐船之脚船，枋楣橡柱如瓜架者谓之「丝瓜架」。木顶船谓之「飞仙」，制如苏州酒船，本于城内沙氏所造，今谓之「沙飞」，皆用篙。沙飞梢舱有灶，无灶者谓之「江船」，用橹者为「摇船」，前席棚后木顶者谓之「牛舌头」，用桨者为「划子船」，双桨为「双飞燕」，亦曰「南京篷」…沙飞重檐飞舻，有小卷棚者谓之「太平船」，覆者为「顶」，以玻璃嵌窗者谓之「玻璃船」。至于四方客卿达官以及城内仕宦向有官船。皆住北门马头，非游人所得乘也。（卷19，第403页)
\end{BookQuote}

可知画舫主要是停靠在岸边用于娱乐的船，这种船已经不再用于水运，经过改造装修，提供水上娱乐场所；画舫也可指湖上的游船。《扬州画舫录》第十一卷记载虹桥与画舫之游的内容很丰富。

扬州的游览胜地当属北郊的瘦西湖，尤以虹桥为标志。李斗说：「游人泛湖，以秋衣、蜡屐打包，茶融、灯遮，点心、酒盏，归之茶担，肩随以出。」（卷11，第240页）若治具待客湖上，先投柬帖，上书「湖舫候玉」。相沿成俗，以此为礼，平时招携游赏，不必投柬。

画舫有市有会。春为梅花、桃花二市，夏为牡丹、芍药、荷花三市，秋为桂

\vfill
\begin{BookNotes}
\BookNote{22.}{冯尔康、常建华，《清人社会生活》，天津：天津人民出版社，1990年版，第474页。}
\end{BookNotes}

% 原书第 242 页（PDF 第 255 页）
\BookPage{242}{255}
\noindent 花、芙蓉二市；又正月财神会市，三月清明市，五月龙船市，六月观音香市，七月盂兰市，九月重阳市。李斗说，每当有市，游人众多，船价数倍。

游船中的「龙船」最为壮观。龙船长十余丈，前为龙首，中为龙腹，后为龙尾，各占一色。四角枋柱，旗帜招展。龙船自五月朔至十八日为一市。先于四月晦日试演，谓之「下水」；至十八日牵船上岸，谓之「送圣」。这时正值端午节，有水上的各种表演。如龙船尾长丈许，牵彩绳令小儿水嬉，谓之「掉梢」。有「独占鳌头」、「红孩儿拜观音」、「指日高升」、「杨妃春睡」诸戏。再如小船载乳鸭，往来画舫间，游人鬻之掷水中，龙船执戈竞斗，谓之「抢标」。端午之后，外河徐宁、缺口诸门龙船由响水闸牵入内河，称为客船。「送圣」后奉太子（屈原）于画舫中礼拜，祈祷收灾降福，据说是「举国若狂」。（卷11，第240-241页）

画肪有堂客、官客之分，堂客为妇女之称。妇女上船，四面垂帘，屏后另设小室如巷，香枣厕筹，位置洁净。船顶皆方，可载女舆。家人挨排于船首，以多为胜，称为堂客船。一年中惟龙船市堂客船最多。迨至灯船夜归，弃舟登岸，天宁寺前，拱宸门外，肩舆等候堂客回家。（卷11，第241页）

扬州夜游生活较为丰富，形成了「夜生活」。李斗说:「城内富贵家好昼眠，每自旦寝，至暮始兴，燃烛治家事，饮食燕乐，达旦而罢，复寝以终日。由是一家之人昼睡夕兴，故泛湖之事，终年不得一日领略。即有船之家，但闲泊浦屿，或偶一出游，多于申后酉初，甫至竹桥，红日落尽，习惯自然。」（卷11，第241页）

不过富贵大家族一旦出游，气势非凡。他们以大船载酒，数艘并集，衔尾以进。盛至三舟并行，宾客喧阗，遥看如驾山倒海而来。

扬州城内画舫无灶，惟「沙飞」之船有之，故多以沙飞代酒船。李斗说，城中奴仆善烹饪者，为家庖；有以烹饪为佣赁者，为外庖。其自称曰厨子，称诸同辈曰厨行。游人赁以野食，乃上沙飞船，令工人肩挑，谓之厨担。厨子随其后，各带所用之物，裹之以布，谓之刀包。于是画舫在前，酒船在后，「橹篙相应，放乎中流，传餐有声，炊烟渐上」，谓之行庖。（卷11，第242页）李斗还说:烹饪之技，家庖最胜，如吴一山炒豆腐，田雁门走炸鸡，江郑堂十样猪头，汪南溪拌


% 原书第 243 页（PDF 第 256 页）
\BookPage{243}{256}
\noindent 鲟鲤，施胖子梨丝炒肉，张四回子全羊，汪银山没骨鱼，江文密车螯饼，管大骨董汤、鱼糊涂，孔言刃庵螃蟹面，文思和尚豆腐，小山和尚马鞍乔，风味皆臻绝胜。

湖中有一种歌船。歌船在座船前，歌船逆行，座船顺行，使船中人得与欣赏歌者。歌以清唱为上，十番鼓次之，还有锣鼓、马上撞、小曲、摊簧、对白、评话之类佐胜。

清唱以笙笛、鼓板、三弦为道具。每一市会，争相斗曲，以画舫停篙就听者多少为胜负。多以熙春台、关帝庙为清唱之地。李斗认为，扬州人好度曲而不佳，近时以叶广平唱口为最，着《纳书楹曲谱》，为世所宗，其余无足数。清唱以外、净、老生为大喉咙，生、旦词曲为小喉咙，丑、末词曲为小大喉咙。扬州刘鲁瞻工小喉咙，为刘派，兼工吹笛。大喉咙以蒋铁琴、沈苕湄二人为最，为蒋、沈二派。蒋本镇江人，居扬州，以北曲胜，沈以南曲胜。此外，扬州还有苏州来的大喉咙。可知清唱的流派很多。李斗还说：「扬州清唱既盛，串客乃兴，王山霭、江鹤亭二家最胜，次之府串班、司串班、引串班、邵伯串班，各占一时之胜。其中刘禄观以小唱入串班为内班老生，叶友松以小班老旦入串班，后得瓜张插花法；陆九观以十番子弟入串班，能从吴暮桥读书，皆其选也。」（卷11，第244页）如此更是清唱百花手艳。

十番鼓者，只用笛、管、箫、弦、提琴、云锣、汤锣、木鱼、檀板、大鼓十种，故名十番鼓。番者更番之谓，有《花信风》、《双鸳鸯》、《风摆荷叶》、《雨打梧桐》诸名。后增星钹，器辄不止十种，遂以星、汤、蒲、大、各、勺、同七字为谱。若夹用锣铙之届，则为粗细十番。

锣鼓盛于上元、中秋二节，以锣鼓铙钹。南巡时延师演习，谓之办差锣鼓。

马上撞」即军乐演唱乱弹戏文，城中市肆剪生、开张及画舫、财神、三圣诸会多用之。

小唱以琵琶、弦子、月琴、檀板合动而歌。李斗说，二十年前尚哀泣之声，谓之「到春来」，又谓之「木兰花」；后以下河土腔唱〈剪靛花〉，谓之「网调」。近来群尚〈满江红〉、〈湘江浪〉，皆本调也。其京舵子、起字调、马头调、南京


% 原书第 244 页（PDF 第 257 页）
\BookPage{244}{257}
\noindent 调之类，传自四方，间亦效之。当时有不少名家，如陈景贤善小曲，兼工琵琶，人称为「飞琵琶」；潘五道士能吹无底洞箫以和小曲，称名工；苏州牟七以小唱冠江北，后多须，人称为牟七胡子；朱三工四弦，江鹤亭招之入康山草堂。

评话盛于江南，如柳敬亭、孔云霄、韩圭湖诸人，屡为陈其年、余淡心，杜茶村、朱竹垞所赏鉴。扬州平话「称绝技者，吴天绪《三国志》，徐广如《东汉》，王德山《水浒记》，高晋公《五美图》，浦天玉《清风闸》，房山年玉蜻蜓，曹天衡《善恶图》，顾进章《靖难故事》，邹必显《飞驼传》，谎陈四《扬州话》，皆独步一时」。（卷11，第246页）大鼓书始于渔鼓简板说孙猴子，佐以单皮鼓檀板，谓之「段儿书」；后增弦子，谓之「靠山调」。

湖上画舫，有不少民间艺人表扬，高超技艺受到游人称赞。如大松、小松兄弟，本浙江世家子，落拓后卖歌虹桥。大松弹月琴，小松拍檀板，就画舫互唱觅食。再如井天章善学百鸟声，游人每置之画舫间与鸟斗鸣，其技与画眉杨并称。还有匡子驾小艇游湖上，以卖水烟为生。以吐烟为图案著称，时人谓之「匡烟」，遂自名其船曰「烟艇」。

画舫多作牙牌、叶格、纸牌诸戏，吸引游客，招揽酒食生意。牙牌以竹代替，四人合局，得四为上。二人对局为「扛」，有苏、扬之分。三四人合局，以点大得者为负，谓之「挤黄」。叶格以「马吊」为上，马吊四十张，自空堂至于万万贯，十万贯以下，均易被攻。纸牌始用三十张，即马吊去十子一门，谓之「斗混江」；后倍为六十，谓之「挤矮」；又倍之为一百二十张。五人斗，人得二十张。近今尽斗十壶，又增以福、禄、寿、财、喜五星，计张一百二十有五。

画舫多以弈为游者，围棋国手聚集扬州。扬州国手只韩学元一人；寓公则有樊麟书、程懒予、周东侯、盛大有、汪汉年、黄龙士、范西屏、何暗公、施定庵、姜起士诸人。施定庵父殁，从母改适范氏，生西屏。施、范同时称国手，范着《桃花泉弈谱》，施着《弈理指归》，皆传于世。当时言棋者，动称施、范。程兰如弈棋不及施、范，而象棋称国手。

游艺活动还有博戏。古博戏跌成，亦称「纯」，时人谓之「拾博」。用三钱者


% 原书第 245 页（PDF 第 258 页）
\BookPage{245}{258}
\noindent 为三星，六钱者为六成，八钱者为八又。「均字均幕为成，四字四幕为天分。天分必幕与幕偶，字与字偶，长一尺，不杂不斜，以此为难。」（卷16，第351页）李斗说：此技遍于湖上，所博之物，以茉莉、玫瑰二花最多。

湖上的花船、灯船也是重要的风景。灯船多用鼓棚，楣枋力檐，以宫灯为最丽。其次琉璃，一船连缀百余而出。或值良辰令节，诸商各于工段临水张灯，两岸中流，交辉焕朵。「花船于市会插花画舫中，大者用磁缸，小则瓶洗之属，一瓶动值千金。插花多意外之态，此技瓜洲张某最优，时人称为瓜张。优者叶友松一人，亦传其法。十番教师朱五呆亦能之。」（卷11，第250页）

湖岸长堤有人贩卖饮食者。乔姥置大茶具，旁列茶盒，矮竹几杌数十。每茶一碗二钱，称为「乔姥茶桌子」。每龙船时，茶客往往不给钱而去。还有卖糖的，有一种糖叫大观楼，以紫竹作担，列糖于上，糖中裹盐脂豆馅之类，贵至十数钱一枚。又有提篮鸣锣唱卖糖官人、糖宝塔、糖龟儿者，味不甚佳，止供小儿。「或置竹钉数十于竹筒中，其端一赤而余皆黑，以钱贯之。适中赤者则得糖，否则负。口中唤唱，音节入古。」（卷11，第250-251页）

清明前后，是出游的黄金季节。肩担卖食者多俊秀少年，竞尚妆饰。「每着蓝藕布衫，反纫钩边，缺其衽，谓之琵琶衿。缝错伍取窄，谓之棋盘裆。草帽插花，蒲鞋染蜡。」（卷11，第251页）卖豆腐脑、茯苓糕，唤声柔雅可听。夏月则有卖洋糖豌豆，秋月有卖芋苗子的普通贩卖者。

北方人王蕙芳以卖果子为业。清晨以大柳器贮各色果子，先货于苏式小饮酒肆，次及各肆，其余则于长堤售尽。自称为「果子王」。「其子八哥儿卖槟榔，一日可得数百钱。」（卷11，第251页）

堤上的休闲者，有的展示自己的爱好与特长。谢身山多奇技，会做飞鸟。李斗说他:「每和泥贯以竹蔑，置数十枚于袖中。行至堤上，启袖放之，如燕雀腾飞，轧轧有声。」（卷11，第251页）

每天清晨，城中养鸟者携白翎雀于堤上学黄鹂声。白翎雀本北方鸟，江南人好之，饲于笼中，一鸟动辄百金。笼之价值，贵者如金戏盆，中铺沙斫石，令雀于其


% 原书第 246 页（PDF 第 259 页）
\BookPage{246}{259}
\noindent 上鼓翅，谓之「打蓬」。李斗记载说：「若画舫中，每悬之于船楣，以此为戏。(卷11，第251页）此外，养画眉等鸟类者也很多。堤上多蝉，人们用长竿黏落，贮以竹筐，沿堤货麦，以供儿童嬉戏。

堤上也可以看到傀儡戏与猴戏的表演。凤阳人精于此道，他们蓄猴令其自为冠带演剧，谓之猴戏。「又围布作房，支以一木，以五指运三寸傀儡，金鼓喧阗，词白则用叫颡子，均一人为之，谓之肩担戏。」（卷11，第251页）李斗说，此二戏正月城内极多，皆预于腊月抵郡城，寓文峰塔壶芦门客舍。至元旦进城，上元后城中表演已遍，出城在堤上表演。

其实，堤上已经形成了公共表演空间。堤上的杂技表演尤能说明此点，李斗详细记载了这一情形：

\begin{BookQuote}
杂耍之技，来自四方，集于堤上。如立竿百仞，建帜于颠，一人盘空拔帜，如猱升木，谓之「竿戏」。长剑直插喉嗉，谓之「饮剑」。广筵长席，灭烛罨火，一口吹之，千碗皆明，谓之「壁上取火，席上反灯」。长绳高系两端，两人各从两端交过，谓之「走索」。取所佩刀令人尽力刺其腹，刀摧腹皤，谓之「弄刀」。置盘竿首，以手擎之，令盘旋转；复两手及两腕、腋、两股及腰与两腿，置竿十余，其转如飞。或飞盘空际，落于原竿之上，谓之「舞盘」。戏车一轮，中坐数女子，持其两头摇之，旋转如环，谓之「风车」。一人两手执箕，踏地而行，扬米去糠，不溢一粒，谓之「簸米」。置丈许木于足下，可以超乘，谓之「踩高𫏋」。以巾覆地上，变化什物，谓之「撮戏法」。以大碗水覆巾下，令隐去，谓之「飞水」。置五红豆于掌上，令其自去，谓之「摘豆」。以钱十枚，呼之成五色，谓之「大变金钱」。取断臂小儿，令吹笙，工尺俱合，谓之「仙人吹笙」。癸丑秋月，诸杂耍醵资买棹，聚于熙春台，各出所长，凡数日而散。一老人年九十许，曳大竹重百余斤，长三四丈，立头上，每画舫过，与一钱。黄文旸为之立传。（卷11，第252页)
\end{BookQuote}

李斗还例举了一些堤上的表演者，如：

\begin{BookQuote}
汪某以串客倾其家，至为乞儿。遂傅粉作小丑状，以五色笺纸书其招牌，
\end{BookQuote}


% 原书第 247 页（PDF 第 260 页）
\BookPage{247}{260}
\begin{BookQuote}
曰「太平一人班」。「有招之者，辄出戏简牌，每出价一钱。」（卷11，
\end{BookQuote}

第252页）

\begin{BookQuote}
有一个叫王大的人，「置碗头上，碗中立纸绢人数寸，跪拜跳跟，至于偃仆，其碗不坠」。（卷11，第252页）后来安庆武部习其技，置灯头上，谓之「滚
\end{BookQuote}

灯。

\begin{BookQuote}
北方人宋二，「好与禽兽伍，禽兽亦乐与之狎」。（卷11，第253页）
\end{BookQuote}

还有摔跤表演。两人裸体相扑，谓之「摆架子」。

江宁人造方圆木匣，中点花树、禽鱼、怪神、秘戏之类，外开圆孔，蒙以五色瑇瑁，一目窥之，谓之「西洋镜」。

这里还贩卖特色率要纪念品。北郊多萤，土人制料丝灯，以线系之，于线孔中纳萤；其式方、圆、六角、八角及画舫、宝塔之属，谓之「火萤虫灯」。「近多以蜡丸𦶟之，每晚揭竿首鬻卖，游人买作土宜。」（卷11，第253页）还有一种「西瓜灯」，取西瓜皮镂刻人物、花卉、虫鱼之戏。城内还多用料丝作大山水灯片。

配合游湖，还有餐饮服务。玉版桥王廷芳茶桌子最着，与双桥卖油糍之康大合本，各用其技。游人至此半饥，茶香饼熟，颇易得钱。

堤上还有迎候游船的「乞丐」表演民歌。李斗说:「玉版桥乞儿二，一乞剪纸为旗，揭竹竿上，作报喜之词；一乞家业素丰，以好小曲荡尽，至于丐，乃作男女相悦之词，为《小郎儿曲》。相与友善，共在堤上。每一船至，先进《小郎儿曲》，曲终继之以报喜，音节如乐之乱章，人艳听之。《小郎儿曲》即《十二月》、《探茶》、《养蚕》诸歌之遗，呢呢儿女语，恩怨相尔汝。词虽鄙俚，义实和平，非如市井中小唱淫靡媚亵可比。」（卷11，第254页）李斗还说，扬州的这些民歌传播到广州，非常有名，珠江「河网船中冒作扬妓者。其唱则以是曲为土音，岭外传之，及于惠、潮，与木鱼、布刀诸曲相埒。郡中剞劂匠多刻诗词戏曲为利，近日是曲翻板数十家，远及荒村僻巷之星货铺，所在皆有，乃知声音之道，感人深也」。（卷11，第254页)


% 原书第 248 页（PDF 第 261 页）
\BookPage{248}{261}
画舫游者的饮食方法比较多元化。李斗说:「野食谓之『饷』。画舫多食于野，有流觞、留饮、醉白园、韩园、青莲社、留步、听箫馆、苏式小饮、郭汉章馆诸肆，而四城游人又多有于城内肆中预订者，谓之『订菜』，每晚则于堤上分送各船。」（卷11，第254-255页）

扬州作为旅游城市，娼妓也是娱乐的一个方面。李斗说，清初官妓，谓之「乐户」。至康熙间，裁乐户，遂无官妓。官妓既革，土娼潜出，如私窠子、半开门之属，有司禁之。泰州有渔网船，「郡城呼之为网船滨，遂相沿呼苏妓为苏滨，土娼为扬滨，一逢禁令，辄生死逃亡不知所之」。（卷9，第189页）李斗介绍道：

\begin{BookQuote}
自龙头至天宁门水关，夹河两岸，除各有可记载者，则详其本末，若夫歌喉清丽、技艺共传者，则不能枚举。如白四娘者，扬州人，因县吏朱某曾拯其急难，后朱缘事几置法。伊倾家谋救，得充边远军，不至死。赵大官、赵九官、大金二官、小金二官、陈银官、巧官、麻油王二官、杨大官、杨三官、吴新官、汪大官、闵得官、闵二官、沈四官、沈大二官、赵三官、陆爱官、佟凤官、夏大官、小青青、蒋大官、蒋二官、张三官、王大官、小脚陈三官、大脚陈三官，此皆色技俱佳，每舟游湖上，遇者皆疑为仙至。若面店王三官者，则又开扬州苏滨之鼻祖者矣。以技艺见重，不以色也。其妾五官娟好，然遇冶游询其年齿姓字，则面红潜遁，此又苏滨中之奇人焉。若高小女子，本系扬人，丰姿绝世，而才艺一时无两，徐九官与之齐名，其实则逊之甚远也。陈巫云、琼子、凤子、南门高二官、李二官、兴化李二官、蒋六子、小丁香、郭三、三扬。陈四俗呼为「盐豆子」，有女梅梅，年十四，真绝色，后为有力者购去，冀北之群空矣。（卷9，第193页)
\end{BookQuote}

由此可见，扬州游船所经两岸，多有「色技俱佳」的娼优，以备游湖客人所需。当时「沿着新城运河两侧发展起一个类似的青楼区域」23。李斗还说，扬州娱乐圈有品评倡优的现象，如顾姬，行四，字霞娱。工词曲，解诗文，住姜家墩天心庵旁。「会钱湘舱三元棨过扬州，于谢未堂司寇公宴席中品题诸妓，以杨小保为女状元，霞娱为女榜眼，杨高为女探花。」（卷9，第202页）扬州的这种灯红酒绿的生活，

\vfill
\begin{BookNotes}
\BookNote{23.}{安东篱，《说扬州：1550-1850年的一座中国城市》，中华书局，2007年版，第161页。}
\end{BookNotes}

% 原书第 249 页（PDF 第 262 页）
\BookPage{249}{262}
\noindent 被认为是身处醉乡:「时际升平，四方安乐，故士大夫俱尚豪华，而尤喜狭邪之游。在江宁则秦淮河上，在苏州则虎丘山塘，在扬州则天宁门外之平山堂，画船箫鼓，殆无虚日。妓之工于一艺者，如琵琶、鼓板、昆曲、小调，莫不童而习之，间亦有能诗画者，能琴棋者，亦不一其人。流连竟日，传播一时，才子佳人，芳声共着。然而以此丧身破家者有之，以此败名误事者有之，而人不知醒，譬诸饮酒，常在醉乡，是诚何心哉！」24

五戏曲欣赏

清代是戏曲发展时期。清初，明末以来盛行的昆山腔与弋阳腔处于正统地位，但是已暴露出日益脱离群众的情形，特别是昆腔更为明显。内容上多是通过帝王将相、才子佳人宣传封建礼教；形式上不外是悲欢离合的套子，艺术僵化。不过流传到民间的昆、弋诸腔同各地的民歌、说唱以及民间小戏结合，向着地方化的方向发展，还有一些是新产生的地方戏。

扬州是地方戏演出繁盛的都市。李斗说：「两淮盐务例蓄花、雅两部，以备大戏：雅部即昆山腔；花部为京腔、秦腔、弋阳腔、梆子腔、啰啰腔、二簧调，统谓之『乱弹」。昆腔之胜，始于商人徐尚志征苏州名优为老徐班；而黄元德、张大安、汪启源、程谦德各有班。洪充实为大洪班，江广达为德音班，复征花部为春台班；自是德音为内江班，春台为外江班。今内江班归洪箴远，外江班隶于罗荣泰。此皆谓之『内班』，所以备演大戏也。」（卷5，第103页）盐商的戏班子众多，除了自己平日欣赏戏曲外，皇帝南巡时的迎驾演出是重要的事情。

两淮盐商凭着自己强大的经济力量和与官府的关系，自养戏班，为皇家服务，亦称作「内班」。戏班也有相当多的则属于民间经营，农村有不少艺人组织戏班在本地演出，如扬州「城外邵伯、宜陵、马家轿、僧道桥、月来集、陈家集人，自集成班」（卷5，125页），属于扬州郊区的土班。艺人的报酬有所不同，一般演员收入不多，仅可糊口，但名角收入颇丰。扬州的苏州艺人在乾隆末年角色优劣，以戏钱多寡为差，有七两三钱、六两四钱、五两二钱、四两八钱、三两六钱之分，内班

\vfill
\begin{BookNotes}
\BookNote{24.}{（清）钱泳《履园丛话》卷7〈臆论·醉乡〉，中华书局，1979年版，上册，第193页。}
\end{BookNotes}

% 原书第 250 页（PDF 第 263 页）
\BookPage{250}{263}
\noindent 角色皆七两三钱。人数之多，至百数十人，「此一时之胜也」。（卷5，117页）

扬州城内苏唱街老郎堂，是梨园总局。每一班入城，先于老郎堂祷祀，谓之「挂牌」；次于司徒庙演唱，谓之「挂衣」。每团班在中元节，散班在竹醉日。梨园以副末开场，为领班；副末以下老生、正生、老外、大面、二面、三面七人，谓之「男脚色」；老旦、正旦、小旦、贴旦四人，谓之「女脚色」；打译一人，谓之「杂」。李斗说:「此江湖十二脚色，元院本旧制也。」（卷5，第117页）

著名戏班派头大，但是有时遇到懂戏者也会吃到苦头。李斗讲了老徐班的两个故事，一则是：

\begin{BookQuote}
纳山胡翁尝入城订老徐班下乡演关神戏，班头以其村人也，绐之曰：「吾此班每日必食火腿及松萝茶，戏价每本非三百金不可。胡公一一允之。班人无已。随之入山，翁故善词曲，尤精于琵琶。于是每日以三百金置戏台上，火腿松萝茶之外，无他物。日演《琵琶记》全部，错一工尺，则翁拍界尺叱之，班人乃大惭。（卷5，第130页）
\end{BookQuote}

另一则是：

\begin{BookQuote}
西乡陈集尝演戏，班人始亦轻之，既而笙中簧坏，吹不能声，甚窘。詹政者，山中隐君子也，闻而笑之，取笙为点之，音响如故，班人乃大骇。詹徐徐言数日所唱曲，某字错，某调乱，群优皆汗下无地。（卷5，第130页)
\end{BookQuote}

这两个故事表明，扬州民间戏曲文化深厚，有戏曲欣赏的行家里手。

戏剧的繁荣，必定会出现从临时场所演出到固定场所演出，从不定时演出到定时演出，从集资（凑份子）请演员演出到在营业性场所看戏过渡的趋势，这就带来了戏园的出现和发展。清代扬州就有营业性的戏园，它是演出方式变化的结果。李斗说：「天宁寺，本官商士民祝厘之地。殿上敬没经坛，殿前盖松棚为戏台，演仙佛、麟凤、太平击壤之剧，谓之大戏。事竣拆卸。迨重宁寺盖大戏台，遂移大戏于此。」（卷5，103页）说明由于演出活动的频繁，已有固定的大戏台取代了临时性的木棚戏台，这个变化发生在乾隆末年。嘉庆十三年（1808），又仿照北京连续设


% 原书第 251 页（PDF 第 264 页）
\BookPage{251}{264}
\noindent 置了三个专门的戏馆，三月新城有总商余晟瑞家闲园出赁为固乐园，闰五月旧城大东门内的胜春园酒肆改为阳春茶社，六月新城兴教寺后岑姓废园改为丰乐园，这些戏馆「踵事增华，聊以待腰缠之集」25。

扬州的戏剧演出种类繁多。李斗说道:

\begin{BookQuote}
郡城花部，皆系土人，谓之「本地乱弹，此土班也。至城外邵伯、宜陵、马家桥、僧道桥、月来集、陈家集人，自集成班，戏文亦间用元人百种，而音节服饰极俚，谓之「草台戏」。此又土班之甚者也。若郡城演唱，皆重昆腔，谓之「堂戏」。本地乱弹只行之祷祀，谓之「台戏」。迨五月昆腔散班，乱弹不散，谓之「火班。后句容有以梆子腔来者，安庆有以二簧调来者，弋阳有以高腔来者，湖广有以罗罗腔来者，始行之城外四乡，继或于暑月入城，谓之「赶火班」。而安庆色艺最优，盖于本地乱弹，故本地乱弹间有聘之入班者。京腔用汤锣不用金锣，秦腔用月琴不用琵琶，京腔本以宜庆、萃庆、集庆为上。（卷5，第125页）
\end{BookQuote}

当地演出的「本地乱弹」使用元人戏文，音节服饰土气，称之为「草台戏」，城里重视的是作为「堂戏」的昆腔，当时士大夫以及富商喜欢的剧种。同时各地的剧种来到扬州，先在城外四乡，继之入城演出，尤以安庆二簧调、京腔受人欢迎。

扬州的「本地乱弹」在于其他地方戏的交流中发生变化。李斗说:「郡城自江鹤亭征本地乱弹，名春台，为外江班。不能自立门户，乃征聘四方名旦，如苏州杨八官、安庆郝天秀之类；而杨、郝复采长生之秦腔，并京腔中之尤者，如《滚楼》、《抱孩子》、《卖》、《送枕头》之类，于是春台班合京秦二腔矣。熊肥子演《大夫小妻打门吃醋》，曲尽闺房儿女之态。」（卷5，第125页）扬州为地方戏的发展提供了舞台。其实，扬州的戏曲演出还为京剧的形成做出了贡献，李斗指出:「自四川魏长生以秦腔入京师，色艺盖于宜庆、萃庆、集庆之上，于是京腔效之，京秦不分。迨长生还四川，高朗亭入京师，以安庆花部，合京秦两腔，名其班曰三庆，而曩之宜庆、萃庆、集庆遂湮没不彰。」（卷5，第125页)三庆、春台，即属于四大徽班，他融合了安庆花部与京、秦两腔，而所谓徽班正是扬州盐商

\vfill
\begin{BookNotes}
\BookNote{25.}{（清）林苏门撰，《邗江三百吟》卷8〈戏馆〉引言，广陵书社，2005年版，第108页。}
\end{BookNotes}

% 原书第 252 页（PDF 第 265 页）
\BookPage{252}{265}
\noindent 支持下的戏班子，扬州地方戏之间的交流，成就了艺人的戏曲艺术。如樊大演《思凡》一出，始则昆腔，继则梆子、罗罗、弋阳、二簧，无腔不备，人称「戏妖」。「陆三官花鼓得传，而熟于京秦两腔。」（卷5，第126页）京师萃庆班谢瑞卿，工《水浒记》之阎婆惜。「关大保演阎婆惜效之，自是扬州有谢氏一派。」（卷5，第126页）四川魏长生以演出秦腔出名，曾来扬州「投江鹤亭，演戏一出，赠以千金。」（卷5，第126页）地方戏曲艺术的相互吸收，还具体表现在：

\begin{BookQuote}
凡花部脚色，以旦丑、跳虫为重，武小生，大花面次之。若外末不分门，统谓之「男脚色」。老旦、正旦不分门，统谓之「女脚色」。丑以科译见长，所扮备极局骗俗态，拙妇骏男，商贾刁赖，楚咻齐语，闻者绝倒。然各囿于土音乡谈，故乱弹致远不及昆腔。惟京师科浑皆官话，故丑以京腔为最。如凌云浦本世家子，工诗善书，而一经傅粉登场，喝采不绝。广东刘八，工文词，好驰马，因赴京兆试，流落京腔，成小丑绝技。此皆余亲见其极盛，而非土班花面之流亚也。吾乡本地乱弹小丑，始于吴朝、万打岔，其后张破头、张三网、痘张二、郑士伦辈皆效之。然终止于土音乡谈，取悦于乡人而已，终不能通官话。近今春台聘刘八入班，本班小丑效之，风气渐改。（卷5，第126-127页)
\end{BookQuote}

可知扬州本地乱弹春台班吸收京腔小丑念白，改为官话。戏曲服装也发生了交融改变。戏具谓之「行头」，行头分衣、盔、杂、把四箱。衣箱中有大衣箱、布衣箱之分。大衣箱、盔箱又细分文扮、武扮、女扮三种类型。「盐务自制戏具，谓之『内班行头』，自老徐班全本《琵琶记》请郎花烛，则用红全堂，风木余恨则用白全堂，备极其盛。他如大张班，《长生殿》用黄全堂，小程班《三国志》用绿虫全堂。小张班十二月花神衣，价至万金；百福班一出北饯，十一条通天犀玉带；小洪班灯戏，点三层牌楼，二十四灯，戏箱各极其盛。若今之大洪、春台两班，则聚众美而大备矣。」（卷5，第129-130页）在扬州地方戏交流中日渐成熟的戏曲艺术，终于进京后大放异彩，促进了京剧的形成。

清代以花部为标志的地方戏十分兴盛，清廷担心戏文有不利于统治的内容。乾隆四十二年（1777）「巡盐御史伊龄阿奉旨于扬州设局改曲剧，历经图思阿并伊公两任，凡四年事竣。总校黄文旸、李经，分校凌廷堪、程枚、陈治、荆汝为，委员


% 原书第 253 页（PDF 第 266 页）
\BookPage{253}{266}
\noindent 淮北分司张辅、经历查建佩、板浦场大使汤惟镜。」（卷5，第103页）在扬州设局检查剧本，也说明扬州戏曲兴盛。乾隆帝指示：剧本中「如明季国初之事，有关涉本朝字句…至南宋与金朝关涉词曲，外间剧本，往往有扮演过当，以致失实者；流传久远，无识之徒，或至转以剧本为真，殊有关系，亦当一体饬查」26。看来对剧本的检查，主要出于民族关系，即把汉人斥责女真、满族的地方加以修改，以免除民人对清朝的正统地位产生怀疑。剧本修改既成，黄文旸着有《曲海》二十卷，其序云:「乾隆辛丑间，奉旨修改古今词曲。予受盐使者聘，得与修改之列，兼总校苏州织造进呈词曲。因得尽阅古今杂剧传奇，阅一年事竣。」（卷5，第107页）当时南方的另一个戏曲演出中心苏州呈交收到的剧本。

李斗推崇著名戏剧家:「程志辂，字载勋，家巨富，好词曲。所录工尺曲谱十数橱，大半为世上不传之本。凡名优至扬，无不争欲识。有生曲不谙工尺者，就而问之。子泽，字丽文，工于诗，而工尺四声之学，尤习其家传。」（卷5，第130页）不过到了李斗生活的时代，「胡翁久没，詹亦下世，惟程载勋尚存，然亦老且贫，曲本亦渐散失。德音班诸工尺，汪损之尝求得录之，不传之调，往往而有也」。（卷5，第130页）似乎扬州的戏剧演出风华已过，不再辉煌。

\BookSubheading{六、文人雅集与文物鉴赏}

扬州的繁华，成为文人雅士聚集之地。李斗说：「扬州为南北之冲，四方贤士大夫无不至此。予见闻所囿，未能遍记。有游迹数至而无专主之家，以虹桥为文酒聚会之地。」（卷10，第230-231页）李斗列举了记忆所及到过扬州学者文人，有梅文鼎、朱彝尊、阎若璩、朱筠、钱大昕、王昶、沈初、袁枚、王鸣盛、金榜、李汪度、卢文弨、邵晋涵、陆师、张书勋、钮玉、周大纶、汪启淑、顾文煊、谈泰、蒋莘、尤荫、陈实孙、程世淳、曹文埴、胡先声、耿蕙、奇丰额、江绍莘、赵廷枢、阮承裕、江嘉理、董洵、文元星、汪坤、巴源绶、洪锡恒等人。

来到扬州的文人，往往以文会友，举行雅集。如汪坤，字元至，号玉屏，旌德人。作诗上佳，广交游，尝于扬州集诗人为会，刻有《吟香馆合稿》。「会中李

\vfill
\begin{BookNotes}
\BookNote{26.}{《清高宗圣训》卷264，乾隆四十五年十一月乙酉条。}
\end{BookNotes}

% 原书第 254 页（PDF 第 267 页）
\BookPage{254}{267}
\noindent 天澄，字九渊，号疲仙，穷而工诗；汪俊，字杰士，号碧峰，工诗画；汤振家，字绍先，号绣谷；萧炳，字永着，号晴岩，吴仁煜，字春陵；李澍，字澍千，号渔庄；张镠，字子贞，号老姜，工诗画；秦昱，字德明，号岑棋；李桐，字于汤，号琴轩；许善，字松圃；胡保泰，字东山，山阴人；叶建侯，字冠伯，号春屏，丹徒人；查善，字楚珍，号大其，海昌人°」（卷10，第239页）

扬州文人的诗会不少。城南古渡桥旁，歙县汪氏得九莲庵地，建别墅曰南园。汪玉枢，字辰垣，号恬斋，歙县人。早岁能诗，山林性成。南园之盛，由恬斋开始。康熙间，王躬符曾于是园征《城南宴集诗》，计三十六人，各赋七言古诗一首，一时称为胜游。玉枢年七十，尝于易松滋抱山堂中作重九会，子五俱工诗。

三月上巳文人有修禊的传统，地方官往往组织地方文人雅集。曾燠为两淮盐运使，修禊南园，古用上巳，燠行始秋。与会者绘于图，凡八人。李斗说曾燠莅扬州，旦接宾客，夕诵文史，觞咏多暇，「至于北郊诸名胜，转运燕游唱和」。（卷7，第165页）卢见曾，字抱孙，号雅雨山人，山东德州人。为两淮转运使时，「爱古好事，延接文士，风流文采」27。他「筑苏亭于使署，日与诗人相酬咏，一时文宴盛于江南。」（卷10，第218页）卢见曾也曾修禊虹桥，其时和禊韵者七千余人，编次得三百余卷。乾隆三十年（1765），扬州北郊建拳石洞天、西园曲水、虹桥揽胜、冶春诗社、长堤春柳、荷浦熏风、碧玉交流、四桥烟雨、春台明月、白塔晴云、三过留踪、蜀冈晚照、万松叠翠、花屿双泉、双峰云栈、山亭野眺、临水红霞、绿稻香来、竹楼小市、平冈艳雪二十景。此后，湖上复增绿杨城郭、香海慈云、梅岭春深、水云胜概四景。署中文宴，尝书之于牙牌，以为侑觞之具，谓之「牙牌二十四景」。据李斗说:「座中皆天下士，而贫而工诗者，无不折节下交。」（卷10，第219-220页）如同学者指出的，地方官的上述行为，在于「努力塑造他的社会身分」，而扬州「那是一个盐业贸易能带来巨大商业财富，通过资助就能够买来有声望的作品或是与学者和画家的友谊的地方」28。

文人到扬州，往往为商人家中塾师、清客，儒与商沟通。徽州人业盐扬州，成

\vfill
\begin{BookNotes}
\BookNote{27.}{赵尔巽等撰，《清史稿》卷128，中华书记标点本第37册，第11107页。}
\BookNote{28.}{梅尔清，《清初扬州文化》，上海：复旦大学出版社，2004年版，第82、83页。}
\end{BookNotes}

% 原书第 255 页（PDF 第 268 页）
\BookPage{255}{268}
\noindent 为富商。著名徽商门下，多有文人。如徽商徐氏门下的文人有：

\begin{BookQuote}
齐召南，字次风，号琼台，晚号思园，天台人。幼称神童，年二十三拔贡，馆于徐氏。雍正举博学鸿词，授庶吉士，纂修《大清一统志》。官检讨，校勘经史，纂修《明鉴纲目》、《大清会典》、《文献通考》，勘定《通礼》，称经
\end{BookQuote}

学儒臣。

\begin{BookQuote}
程瑶田，字易田，歙县人。孝廉，官太仓州嘉定县教谕。能诗文，与戴东原、方希园穷经数十年，着《通艺录辨》、《九谷沟洫》诸考，皆能发古人所未发。往来扬州主徐氏。子培，字伯厚，能读父书。今业鹾，寓于扬。
\end{BookQuote}

\begin{BookQuote}
叶敬，字义方，扬州人。工书，馆徐氏，着有《日华堂诗集》。
\end{BookQuote}

\begin{BookQuote}
金农，字寿门，号冬心，浙江仁和人。好古力学，工诗文，精鉴赏，善别古书画，书法汉隶。年五十，妻亡，侨居扬州，从事于画，涉古即古，脱尽画家之习。与徐氏往来，以其学圃改名「交翠林」。
\end{BookQuote}

\begin{BookQuote}
方辅，字密庵，徽州人。工诗，书法苏、米，能擘窠大书。善制墨，来扬州主徐氏，着有诗文集，刻以行世。洪尔鉴，字照堂，歙县人。徐氏姻娅，工诗。
\end{BookQuote}

\begin{BookQuote}
黄占济，字槎山，歙县人。主徐氏，工诗，书法各体备具，为人慷爽。
\end{BookQuote}

\begin{BookQuote}
韩奕，字仙李，扬州人。买园湖上，名曰「韩园」。工诗，善鼓板。蓄砂壶，
\end{BookQuote}

为徐氏客。

\begin{BookQuote}
韩姓，逸其名，扬州人。性古直，为徐氏客。晚年居园后鬼神坛旁，草屋三楹而已，年九十余乃卒。
\end{BookQuote}

文人借助富商谋生，从事文化事业，富商在资助文人与文化事业中提升社会地位。

商人致富，不少人收藏书画古玩，形成风气，文人则是鉴赏家。李斗说「扬州书画家极多，兼之过客往来，代不乏人」。（卷2，第36页）他认为考之志乘，而清朝画家不过十四人，皆扬州土着，书家不过三人，则志乘遗失甚多。扬州的书画家，不少曾馆于各工商家。黄其林，字仰岑，歙县人，寓居江都。子二，长德煦，


% 原书第 256 页（PDF 第 269 页）
\BookPage{256}{269}
\noindent 字次和，「善鉴别古物，所蓄宋元人书画最多」。（卷11，第277页）到过扬州的钱泳说:「吴杜村先生名绍浣，安徽歙县人。世以盐䇲为业，寓扬州已百余年，家道殷富。乾隆乙未、戊戌两科，先生与其兄绍烁同中进士，入翰林。先生精于赏鉴，所藏法书名画甚多。当时如彭南昌、董富阳、王韩城、刘长沙诸相国所贡图籍书画，必经先生品题而后奏进。家有颜鲁公《竹山联句》，徐季海、朱巨川告，怀素小草《千文》，王摩诘《辋川图》，贯休《十八应真像》，皆世间稀有之宝。至宋、元、明人，其次焉者也。」29这是大盐商收藏书画的凸出例证。钱泳《履园丛话》卷二〈阅古〉还列举了众多的扬州所藏文物，如周散邑铜盘，旧藏扬州徐氏，后归洪氏。汉陶陵鼎，为扬州阮云台宫保所藏，阮氏还有汉五凤砖。扬州罗两峰有一汉亭长砖。扬州西山是收藏家喜欢居住的地方，据说:「雍正、乾隆间，镇中号全盛。收藏书画名迹极富，名流多往来于此。」30

不少人艺术上有特长，又精通文物，善于鉴赏。丁敬，字敬身，号钝丁，浙江钱塘人。布衣，酿酒为生。好金石文，穷岩绝壁，手自摹拓，着《武林金石录》。秦汉铜器，宋元名迹，入手即辨。收古泉，多异品。长于诗。「其铁笔有求之者，白十金，为镌一字。」（卷4，第89页）钱苍佩，湖州乌程人，精别宋椠元板。褚竣，字千峰，陕西邰阳人。「以鬻碑版为业，天下金石，搜罗殆尽。」（卷4，第91页）黄德煦，字次禾，博学善识古器。「于古人书画，尤精鉴别。」（卷6，第144页）童岳荐，字砚北，绍兴人，精于盐荚，善谋画，多奇中。童玨，字二树，来扬州时，主其家。「王邃学工诗，善画梅，所藏古今人诗文集殆备，能精别古画、铜磁、玉器、金石、钱刀，足迹遍天下，以所蓄玩好自随。」（卷9，第186页）「武生吴仕柏善鼓琴，日与徐锦堂、沈江门、吴重光、僧宝月游，夜则操缦，三更弗缀。」（卷9，第197页）扬州琴学，以徐炜为最，刊《澄鉴堂琴谱》。次之徐锦堂，着有《五知斋琴谱》，谓之二徐。扬州收藏家多古琴，其最古者，惟罗半查家雷琴，内斫「开元二年雷霄斫」。程邃，字穆倩，号江东布衣，又号垢道人，歙县人。「博学，工诗文，精金石、篆刻、鉴别古书画及铜玉器，家藏亦伙。」（卷10，第214页）善画山水，纯用枯笔，写巨然法，别具神味。工隶

\vfill
\begin{BookNotes}
\BookNote{29.}{（清）钱泳，《履园丛话》卷6〈耆旧·杜邨观察〉，中华书局，1979年版，上册，第157页。}
\BookNote{30.}{（清）林溥撰，《扬州西山小志》之〈轶事十二首〉，广陵书社，2005年版，第22页。}
\end{BookNotes}

% 原书第 257 页（PDF 第 270 页）
\BookPage{257}{270}
\noindent 书，晚年侨居江都。顾廉，字又简，苏州人。精鉴识古器。

扬州的能工巧匠很多，有关工艺品著名。如漆器，「夏漆工娶梨园姚二官之妹为妇，家于头巷，结河房三间。漆工善古漆器，有剔红、塡漆两种，以金银铁木为胎，朱漆三十六次，镂以细锦，合有蔗段、蒸饼、河西、三撞、两撞诸式，盘有方、圆、八角、条环、四角牡丹花瓣诸式；匣有长方两三撞诸式，呼为雕漆器，以此至富，故河房中器皿半剔红，并饰之楯槛，为小秦淮第一朱栏」。（卷9，第189页）

如装潢，吴县叶御夫开设装潢店，「御夫得唐熟纸法，旧画绢地虽极损至千百片，一入叶手，遂为完物。然性孤直，慎结纳，不以技轻许人」。（卷9，第198页）

如刻印琢砚，太仓州人吴献可「通经史，究名法之学，方伯延于家二十年。子完夫，工镌印，琢砚，极奇巧之技」。（卷11，第265页）

如装修，谷丽成精宫室之制，凡内府装修由两淮制造者，图样尺寸，皆出其手。潘承烈，字蔚谷，亦精宫室装修之制。文起，字鸿举，江都人。「精于工程做法，所见古器极多，称赏鉴家。」（卷11，第269页）

\BookSubheading{结语}

综上所述，《扬州画舫录》为我们呈现了盛清扬州繁华的城市生活。我们看到的历史，也是《扬州画舫录》的作者李斗想要告诉我们的历史。

\begin{BookQuote}
《扬州画舫录》属于什么体裁，著名文人袁枚的序将其比之于昔年的杨衒之
\end{BookQuote}

《洛阳名园记》、孟元老《东京梦华录》，属于「标举名胜」之书。而著名官僚兼学者的阮元序说「杨、孟追述往事，此录则目睹升平」。又认为该书仿《水经注》之例，分其地而载之，补郡县志所未载。虽然采及琐事俗谈，属于史家与小说家相通，深合古书体例。

李斗的〈自序〉则说此书在补充志书所略，所以该书「以地为经，以人物记


% 原书第 258 页（PDF 第 271 页）
\BookPage{258}{271}
\noindent 事为纬，按扬州郡城之地…依此次叙之为卷帙，其工段营造之制及画舫之名附于卷末」。李斗说他「幼失学，疏于经史，而好游山水」。自乾隆二十九年至六十年（1764一1795）三十年间，「退而家居，则时泛舟湖上，往来诸工段间」。阅历即熟，以目之所见、耳之所闻加以记录，删而成帙。所以，该书兼备游记与琐闻笔记的性质。因该书具有游记的属性，难怪人们会认为它是旅游指南一类的书。

由于李斗「目睹升平」，该书的政治倾向性也是明显的，阮元序就自称「序其略，俾知吾乡承国家重熙累洽之恩，始能臻此盛也」。

扬州的繁荣，尽管有交通之便，鱼盐之利，为十八世纪盛清的城市繁荣提供了基础，但是康熙、乾隆两位皇帝的南巡驻跸，则提供了特殊的历史机缘。乾隆五十八年（1793）袁枚为《扬州画舫录》所作序言称:「自辛未岁（按:乾隆十六年，1751）天子南巡，官吏因商民子来之意，赋工属役，增荣饰观，由来久矣。即乾隆时期，扬州地方官为了皇帝的巡游，进行以园林为标志的城市建设，使得扬州城市面貌大为改观。阮元序也有类似的表达，说扬州土沃风淳，乾隆帝六次南巡「恩泽稠叠」，于是「士日以文，民日以富」。民富士文成为扬州的新貌。节性老人的跋语也说:「扬州全盛，在乾隆四五十年间…方翠华南幸，楼台画舫，十里不断。」他还说到了嘉庆八年（1803）:「此后渐衰，楼台倾毁，花木凋零。」而道光年间:「闻荒芜更盛。且扬州以盐为业，而造园旧商家多歇业贫散，书馆寒士亦多清苦，吏仆佣贩皆不能糊其口…李艾塘斗撰《画舫录》在乾隆六十年，备载当年景物之盛。」就是说李斗记载了由于乾隆帝南巡致使扬州最为繁华时的盛清城市风貌，嘉道时期扬州逐渐失去昔日的辉煌，走向衰落。

因此，盛清扬州的城市繁华，使得我们在传统中国自然经济下农村、农民的原始封闭的纯朴生活图景之外，看到了商品经济下城市、市民特别是富商、绅士奢华生活的画卷。尽管扬州的繁荣有其商品经济大发展作为支持，然而更重要的是乾隆时期皇帝喜欢在扬州驻跸，留恋这里的山水园亭、戏剧文物，于是官员与富商拨款、投资，改造城市，带来扬州景观的美化、商贸的繁荣，文士的聚集，消费的炫耀。一旦随着清朝盛极而衰，扬州也风华难再。在一定程度上来说，《扬州画舫录》笔下的城市影像，多少因渲染而显得有些虚幻。
